\documentclass[12pt, letterpaper]{article}
\usepackage{amsthm, amsmath, amssymb, graphicx, hyperref, url, pdfpages}
\usepackage[utf8]{inputenc}
\usepackage{pgfplots}
\pgfplotsset{compat=1.18}
\usepackage{tikz}
\usetikzlibrary{arrows.meta, positioning}

\newtheorem{thm}{Theorem}[section]
\newtheorem{lemma}[thm]{Lemma}

\theoremstyle{definition}
\newtheorem{ex}[thm]{Example}
\newtheorem{defn}[thm]{Definition}

\title{RSA Cryptography}
\author{Gordon Chen}
\date{Fall 2026}

\begin{document}
\maketitle

\section{Introduction} \label{intro}
In this paper we will describe RSA cryptography.
In Section \ref{sec_prelim} we will state some results from number theory that are
critical for RSA's correctness.
In Section \ref{sec_rsa} we will describe the RSA cryptography scheme.
In Section \ref{sec_correctness} we will prove that RSA is correct.
In Section \ref{sec_security} we will discuss the security of RSA.
Finally, in Section \ref{sec_history} we will discuss the history of RSA.

\section{Preliminaries} \label{sec_prelim}
First we state some preliminaries that will be critical to our analysis of RSA's correctness later.

\begin{defn}
  We define the integers modulo $n$ as the set of equivalence classes
  $\mathbb{Z}_n = \{[0], [1], ..., [n-1]\}$
  representing the possible remainders after dividing integers by $n$.
\end{defn}

For example, $\mathbb{Z}_3 = \{[0], [1], [2]\}$. For the sake of convenience,
we can also write $\mathbb{Z}_3 = \{0, 1, 2\}$.

We leave it as an exercise to the reader that addition and multiplication are
independent of the choice of representatives in $\mathbb{Z}_n$.
We then define exponentiation modulo $n$ as just repeated multiplication, so
$a^k = \prod_{i=1}^k a \pmod n$.

We say that $a \in \mathbb{Z}_n$ is invertible modulo $n$ if $a$ has
a multiplicative inverse modulo $n$, i.e. if there exists some $a^{-1} \in \mathbb{Z}_n$
such that $a a^{-1} = 1 \pmod n$.

\begin{thm}
  An element $a \in \mathbb{Z}_n$ has a multiplicative inverse if and only if $\gcd(a, n) = 1$.

  \begin{proof}
    Assume that $\gcd(a, n) = 1$.
    By the Bezout Identity, there exists $x,y \in \mathbb{Z}$ such that
    $ax+ny=\gcd(a,n)$, where
    $x$ and $y$ can be found using the extended Euclidean algorithm.
    We leave it as an exercise to the reader to show that Bezout's Identity is true.
    If $\gcd(a,n) = 1$, then we can find $x,y \in \mathbb{Z}$ so
    that $x a + y n = 1$, so $x a = 1 \pmod n$.
    Thus, $x$ is the multiplicative inverse of $a$ modulo $n$.

    Conversely, if $ax = 1\pmod n$, then $ax+ny=1$ for some integer $y$.
    If a common divisor $d$ divides both $a$ and $n$, then $d$ must divide $1$, so
    $\gcd(a, n) = 1$ \cite{conrad_mod}.
  \end{proof}
\end{thm}

We are now ready to state and prove Fermat's little Theorem.

\begin{thm}
  Let $p$ be prime. Then $a^{p-1} = 1 \pmod p$ for $\gcd(a, p) = 1$.

  \begin{proof}
    We first show that $f : \mathbb{Z}_p \to \mathbb{Z}_p$ where $f(n) = a n \pmod p$
    is a bijection.

    We show that $f$ is injective. Let $f(n) = f(m)$. Then $a n = a m \pmod p$.
    Since $\gcd(a, p) = 1$, then $a$ is invertible modulo $p$, so $n = m \pmod p$
    and $f$ is injective.

    Since $\mathbb{Z}_p = \{0, 1, ..., p-1\}$ has $p$ elements, it is a finite set,
    so any injection from $\mathbb{Z}_p$ to itself is a bijection.

    Since $f$ is a bijection, $\mathbb{Z}_p$ and $f(\mathbb{Z}_p)$ are the same sets.
    Then,
    \begin{align*}
      (1)(2)...(p-1) &= (1 a) (2 a) ... (p-1) a \pmod p \\
      (1)(2)...(p-1) &= a^{p-1} (1)(2)...(p-1) \pmod p.
    \end{align*}
    Since $p$ is prime, all numbers $1 \le n < p$ have $\gcd(n, p) = 1$ and are therefore
    invertible modulo $p$. Then right-multiplying by their inverses, we get that
    \begin{align*}
      (1)(2)...(p-1) &= a^{p-1} (1)(2)...(p-1) \pmod p \\
      1 &= a^{p-1} \pmod p
    \end{align*}
    and we are done \cite{conrad_flt}.
  \end{proof}
\end{thm}

We are now ready to talk about RSA.


\section{RSA} \label{sec_rsa}
Here we present the classical RSA cryptography scheme developed in the original RSA paper
which uses Euler's totient function $\phi(N) = (p-1)(q-1)$.
However, note that modern RSA specifications only require $ed = 1 \pmod{\lambda(N)}$, where
Carmichael's function is $\lambda(N) = \text{lcm}(p-1, q-1)$ \cite{rsa}.

Choose (extremely large and random) distinct primes $p$ and $q$ and let $N = pq$.
Let $\phi(N) = (p-1)(q-1)$. Really this is Euler's totient function $\phi$ evaluated at $N = pq$,
but for our purposes it is enough to just take it as a symbol representing $(p-1)(q-1)$.
Then choose a positive integer $e$ with $\gcd(e, \phi(N)) = 1$, and a positive integer $d$
such that $d = e^{-1} \pmod{\phi(N)}$.

The public key is $e, N$ and the private key is $d$.

To encrypt a message $m \in \mathbb{Z}_N$, compute $c = m^e \pmod N$.
To decrypt a piece of ciphertext $c \in \mathbb{Z}_N$, compute $m' = c^d \pmod N$.


\section{Correctness} \label{sec_correctness}
Now we prove RSA's correctness, that decrypting an encrypted message
recovers the original message.

\begin{thm}
  Let $p$, $q$ be distinct primes. Let $N = pq$ and let $e, d$ be positive integers
  such that $d = e^{-1} \pmod{\phi(N)}$.
  Let $m \in \mathbb{Z}_N$. Then $m = (m^e)^d \pmod N$.

  \begin{proof}
    Since $N = pq$, we have that
    \begin{align*}
      m = (m^e)^d \pmod N &\iff m = (m^e)^d \pmod{pq} \\
      &\iff pq | (m^{ed} - m).
    \end{align*}
    Since $p$ and $q$ are distinct primes, then 
    \begin{align*}
      pq | (m^{ed} - m) &\iff p | (m^{ed} - m) \text{ and } q | (m^{ed} - m) \\
                        &\iff m^{ed} = m \pmod p \text{ and } m^{ed} = m \pmod q.
    \end{align*}
    Since $d = e^{-1} \pmod{\phi(N)}$, then $ed = 1 \pmod{(p-1)(q-1)}$, so $ed = k(p-1)(q-1) + 1$
    for some nonnegative integer $k$. Then,
    \begin{align*}
      m^{ed} = m \pmod p &\iff m^{k(p-1)(q-1) + 1} = m \pmod p \\
                        &\iff m \cdot (m^{p-1})^{k(q-1)} = m \pmod p.
    \end{align*}
    If $p | m$ then $m = 0 \pmod p$ and we have that $m^{ed} = 0 = m \pmod p$.

    If $p \nmid m$ then $m$ is invertible modulo $p$, so by Fermat's little Theorem,
    we have that $m^{p-1} = 1 \pmod p$, so $m \cdot (m^{p-1})^{k(q-1)} = m \pmod p$.

    The analysis is exactly the same for $q$. Thus we have shown that $m = (m^e)^d \pmod N$,
    so RSA encryption then decryption recovers the original message \cite{conrad_rsa}.
  \end{proof}
\end{thm}


\section{Security} \label{sec_security}
In the previous section we proved that RSA is correct, that if you encrypt then decrypt a message
using RSA, you will recover the original message.
However, this does not in any way mean that RSA is secure, that adversaries
without the secret key will not be able to decipher the encrypted message.

In RSA, the public key is $e, N$ and the secret key is $d$ where $d = e^{-1} \pmod{\phi(N)}$
and $N = pq$ for primes $p, q$.
If an adversary could find $\phi(N)$, then they could solve for $d = e^{-1} \pmod{\phi(N)}$ by
the extended Euclidean algorithm, and then break encryption since they now know the private key.
Given $N$, if an adversary could factor $N$ and recover $p, q$ such that $N = pq$, then
they could compute $\phi(N) = (p-1)(q-1)$ and then solve for private key $d$.
Thus, it is necessary (but not sufficient, as there may be other attacks)
for RSA security that factoring the product of two large primes is difficult.
There is no proof that factoring, or finding $d$ from $e$ and $N$ is ``hard'', meaning
that there does not exist a polynomial time (classical) algorithm to find factors of
$N$ or recover $d$.
But since currently there is no known (classical) polynomial time algorithm for these 2 problems,
RSA is essentially assumed to be secure, with an important caveat that currently we have only
described textbook RSA, which is deterministic and therefore not a secure encryption scheme since
adversaries can simply encrypt candidate plaintexts and compare ciphertexts. For actual RSA
security, one must further incorporate a randomized encoding scheme, such as RSA-OAEP \cite{oaep}.

As of writing (Fall 2026), the minimum recommended RSA key length is 2048 bits,
meaning that the modulus $N$ should be 2048 bits long, so $p$ and $q$ should be
1024 bits long \cite{rsa}.

Conventionally, $e = 65537 = 2^{2^4} + 1$ is fixed \cite{rsa}. The number $65537$ is
the largest known Fermat prime, a prime of the form $2^{2^n} + 1$ \cite{fermat_prime}.
Its special form allows for efficient modular exponentiation, since computing
$m^{2^{16}+1} \pmod N$ only requires 16 modular squarings to find $m^{2^{16}} \pmod N$
and then a single modular multiplication to find $m^{2^{16} + 1} \pmod N$.
In general, choosing $e = \sum_{i=0}^k b_i 2^i$, where $b_i \in \{0,1\}$ and $b_k \ne 0$, uses $k$
modular squarings and $\sum_{i=0}^k \textbf{1}\{b_i = 1\} - 1$ modular multiplications
with the standard binary exponentiation algorithm.
The choice $e = 65537$ is not always valid for primes $p, q$ since we need $e$ to be invertible
modulo $\phi(N) = (p-1)(q-1)$. If $\gcd(e, \phi(N)) = \gcd(e, (p-1)(q-1)) > 1$, then since
$e = 65537$ is prime, $e|(p-1)$ or $e|(q-1)$, so we just have to regenerate each prime that is congruent
to $1$ modulo $e$ \cite{e_65537}.

In 1991, RSA Laboratories put out the RSA Factoring challenge, a list of RSA keys (moduli)
to factor ranging from 330 bits (100 decimal digits) to 2048 bits (617 decimal digits)
\cite{rsa_challenge}.
The largest of these keys confirmed to be factored is RSA-896, with 896 bits and
270 decimal digits, which was factored on September 19th, 2026 by Stephen A. Weis with the
help of Claude. Factoring RSA-896 used a maximum of 2048 GPUs for a total of
30 GPU-years over 10 days \cite{rsa_896}.
For context, RSA-260 was factored by Eric Lu with the help of Devin at Cognition AI on
September 3rd, 2026, spending 13.5 GPU-years and costing around \$400,000 at market compute prices
\cite{rsa_260}.

While there is currently no known classical algorithm for factoring numbers in polynomial time,
sufficiently large, fault-tolerant quantum computers could do so by using Shor's algorithm \cite{shor}, which poses a threat
for RSA since factoring would no longer be hard \cite{rsa}.

There are many other considerations important for RSA security, including
the necessity of strong pseudo-random number generators, side channel and timing attacks,
as well as robustness against heuristic factoring algorithms \cite{rsa}.


\section{History} \label{sec_history}
In this section, we will give a brief history of RSA.

The RSA cryptosystem is named after Rivest, Shamir, and Adleman, the three inventors of RSA.
Rivest, Shamir, and Adleman invented RSA at MIT in 1977. Interestingly enough, in 1973,
Clifford Cocks invented an equivalent encryption scheme at Britain's Government Communications
Headquarters (GCHQ) intelligence agency. However, Cocks' system was not declassified until 1997
\cite{rsa}.

Today, RSA is widely used for Public Key Infrastructure (PKI) certificates, SSH authentication,
and more, but many systems are migrating towards elliptic curve based cryptosystems
\cite{elliptic_curve} and post-quantum secure lattice based cryptosystems \cite{lattice}.


\clearpage
\newpage
\bibliographystyle{plain}
\bibliography{refs}

% \includepdf[pages=-]{rsa_feedback.pdf}

\end{document}
